% !TeX root = ../papers/nc-p3.tex
\section{負付値座標の境界の全傾きと対角線の近傍}
\label{sec:rank-fifteen-negative-x-boundary}
\begin{theorem}\label{thm:rank-fifteen-negative-x-boundary}
剰余標数3の付値で $\nu(x)<0$, $\nu(y)=2\nu(x)$ を満たす通常な
階数15窓は、剰余傾き $\lambda=\overline{y/x^2}$ の全水準で
標準Koszul左右AS正則であり、各始点は四中心四次式上自由階数256である。
\end{theorem}
\begin{proof}
前節の楕円表示 $W^2=t-t^3$ を使う。\path{computations/verify_rank_fifteen_negative_x_boundary.py}
は $\mathbb F_3(t)$ 上の基底 $1,i,W,iW$ でテンソルを構成する。
四次中心から文字0の二つと文字1の二つを選び、全四始点の十三次までの商を計算する。
逆にした元の分子のノルムと入力分母を除けば、同じ単元主元による計算が特殊化できる。
この平方自由な主元多項式の次数は 118054 である。

$N=4(t-t^3)^2$, $D=(1+t^2)^4$ とすれば $\lambda=N/D$ である。
$0,1,\infty$ の逆像を除いた後、$t\mapsto-t,\ 1/t,\ (1+t)/(1-t)$ の
引き戻しとの最大公約元を安定するまで取る。この八次写像を保つ八変換は
非分岐ファイバー上で推移的である。安定した多項式は次数 840 であり、
有限剰余体で各既約因子の $N/D$ の最小多項式を計算する。
その $N/D$ による引き戻しが、対応する既約因子の積に一致することも確認する。
残った閉点の次数 $1,2,3,4,5,6,8,12$ に対する個数はそれぞれ
$1,2,6,4,2,6,1,1$ である。全23閉点について各一根で平方根を取り、
経路・四次中心と全四始点の中心商を十三次まで再計算する。
この検証はFrobenius共役に及ぶ。分岐水準 $\lambda=1$ は前節で扱った。
各傾きで選んだ $g$ は前節の単元ヤコビ行列で所定の $x,y$ に持ち上がる。
楕円表示で選んだ平方根の符号を取り、中心商の上界を持ち上げると結論が得られる。
\end{proof}

\begin{proposition}\label{prop:rank-fifteen-negative-x-diagonal}
同じ負付値座標の $w,g$ 表示で $\nu(w)>0$ とする。
$\nu(g-w)\ge21\nu(w)$ なら上と同じ正則性と自由性を持つ。
\end{proposition}
\begin{proof}
$w=g=t$ として $\mathbb Q(i)[[t]]$ の二十一桁で計算する。
$A'=(1-s^2)/(2i+w)$, $B'=2s/(2i+w)$, $s=i+w$ とすれば
$n=gA'$, $H/n=gB'$ である。$K$ の定数項は1、$P_\pm$ は2、
$Q_\pm/g$ は $1\mp i$ として平方根を取る。
3の上で単元となる主元を用い、$t$ で割るときは実部・虚部の定数係数がともに
標数零で消えることを確認する。未切断経路の損失は $(0,1,1,3,2)$、
四次中心は $(0,1,2,2)$、五次の積は二桁を使う。
四中心商の六次から十三次の損失は $(4,1,1,1,1,0,1,0)$ となる。
$g=w+w^{21}u$ で $u$ が整なら、単純根の選択により経路係数は
対角線のものと $(w^{21})$ を法として一致する。全消去を通じて一桁が残るため、
この摂動でも中心の持上げと標準商の上界が得られる。
\end{proof}

$g=0$ の一般の辺では $r=gw$ が $g$ の単元倍で、$s=i+w$ は一般の円錐曲線座標である。
従って既知の非分岐極の指数を使える。$w=0$ の辺では上の中心商計算が示すように
一般点で未切断経路の全指数は零である。一方 $w=g$ の未切断四次の16行列のうち
8行列で、指数列は $g=0$ の辺と異なる。このため、この二辺だけで
二変数のSmith指数の補題を適用するための項別等式は成り立たない。
細分の候補 $g=mw$ で $\nu(m)=0$ なら $\nu(y)=\nu(x)$、
$\overline{y/x}=\bar m^4/4$ である。実際 $xw^4\equiv-4$,
$H/w^2\equiv-i\bar m^2$ と前節の $y/x^2$ の式を用いればよい。
